% !TEX root =../ECE343_Slides_Rogers.tex


%%%%%%%%%%%%%%%%%%%%%%%%%%%%%%%%%%%%%%%%%%%%%%%%%%%%%%
%%%%%%%%%%%%%%%%%%%%%%%%%%%%%%%%%%%%%%%%%%%%%%%%%%%%%%

\section{Schedule}
\begin{frame}{}
\renewcommand{\arraystretch}{1}
	\begin{textblock}{14}(0.5,1.75)
	\def\myfontsize{footnotesize}
	\setcounter{mylnum}{24}
%\def\myfontsize{normalsize}
\begin{center}
\begin{\myfontsize}
\begin{tabular}{| c | l | c | l |}
\hline
\textbf{Lesson} 		  	& \textbf{Topics}             & \pbox{5cm}{\textbf{Reading}}	& \textbf{Notes}		\\ \hline\hline
\ifnum\lessonNum=\value{mylnum} \rowcolor{green!20!white}\fi
	\themylnum  
	& BVP (Capacitor, Coax), Electric Potential	
	& 5.16
	&  LO 5.3  
\stepcounter{mylnum}\\ \hline \hline
\ifnum\lessonNum=\value{mylnum} \rowcolor{green!20!white}\fi
	\themylnum 	
	& Electric Boundary Conditions
	& 5.18-5.22	
	& LO 5.4			
\stepcounter{mylnum}\\ \hline
\ifnum\lessonNum=\value{mylnum} \rowcolor{green!20!white}\fi
	\themylnum 	
	& BC's/Capacitance			
	& 5.23-5.24
	& 	
\stepcounter{mylnum}\\ \hline
\ifnum\lessonNum=\value{mylnum} \rowcolor{green!20!white}\fi
	\themylnum 	
	& Electrostatic Lab			
	& Lab packet
	& 	
\stepcounter{mylnum}\\ \hline
\rowcolor{gray!20!white}\multicolumn{4}{| c |}{\textbf{Chapter 6: Steady Current \& Conductivity}}	\\ \hline\hline
\ifnum\lessonNum=\value{mylnum} \rowcolor{green!20!white}\fi
	\themylnum 	
	& Current density, Conductivity, Resistance
	& Ch. 6
	&   					
\stepcounter{mylnum}\\ \hline\hline
\rowcolor{gray!20!white}\multicolumn{4}{| c |}{\textbf{Chapter 7: Magnetostatics}}						\\ \hline\hline
\ifnum\lessonNum=\value{mylnum} \rowcolor{green!20!white}\fi
	\themylnum 	
	&  Magnetic Fields, Gauss' Law for Magnetism	
	& 7.1-7.3		
	&LO 7.1-7.2 \textbf{GR 2} 				
\stepcounter{mylnum}\\ \hline
\ifnum\lessonNum=\value{mylnum} \rowcolor{green!20!white}\fi
	\themylnum 	
	& Ampere's Law 	
	& 7.4-7.5, 7.9
	& LO 7.3		
\stepcounter{mylnum}\\ \hline 
\ifnum\lessonNum=\value{mylnum} \rowcolor{green!20!white}\fi
	\themylnum 	
	&  Ampere's Law Applied
	&  7.6-7.8
	& LO 7.3
\stepcounter{mylnum}\\ \hline 
\ifnum\lessonNum=\value{mylnum} \rowcolor{green!20!white}\fi
	\themylnum	
	& Magnetic Boundary Conditions	
	& 7.10-7.14	
	& LO 7.4 		
\stepcounter{mylnum}\\ \hline 
\ifnum\lessonNum=\value{mylnum} \rowcolor{green!20!white}\fi
	\themylnum  
	& Magnetostatics Lab
	& Lab packet
	&  		
\stepcounter{mylnum}\\ \hline \hline
\rowcolor{gray!20!white}\multicolumn{4}{| c |}{\textbf{Chapter 8: Time-Varying Fields}}					\\ \hline\hline
\ifnum\lessonNum=\value{mylnum} \rowcolor{green!20!white}\fi
	\themylnum 	
	& Faraday's Law, Displacement Current, Maxwell Eqns
	& 8.3, 8.8, 8.9
	& LO 8.1
\stepcounter{mylnum}\\ \hline 
\ifnum\lessonNum=\value{mylnum} \rowcolor{green!20!white}\fi
	\themylnum	
	& Faraday's Law, Displacement Current, Maxwell Eqns 
	&  8.3, 8.8, 8.9
	& LO 8.2-8.3
\stepcounter{mylnum}\\ \hline 

\rowcolor{gray!20!white}\multicolumn{4}{| c |}{\textbf{Chapter 9: Plane Waves in Lossless Media}}					\\ \hline\hline
\ifnum\lessonNum=\value{mylnum} \rowcolor{green!20!white}\fi
	\themylnum	
	& Maxwell's Equations in lossless, source-free region		
	& 9.1-9.2
	& LO 9.1	 
\stepcounter{mylnum}\\ \hline
\ifnum\lessonNum=\value{mylnum} \rowcolor{green!20!white}\fi
	\themylnum	
	& Maxwell's Equations in lossless, source-free region		
	& 9.1-9.2
	& LO 9.1	 
\stepcounter{mylnum}\\ \hline
%\ifnum\lessonNum=\value{mylnum} \rowcolor{green!20!white}\fi
%	\themylnum	
%	& Flex day
%	& 
%	& \textbf{GR 3}			 					
%\stepcounter{mylnum}\\ \hline 
\ifnum\lessonNum=\value{mylnum} \rowcolor{green!20!white}\fi
	\themylnum	
	& Uniform Plane Waves \& Polarization
	& 9.4-9.6
	& LO 9.2-9.3			
\stepcounter{mylnum}\\ \hline 
\ifnum\lessonNum=\value{mylnum} \rowcolor{green!20!white}\fi
	\themylnum	
	& Normal Incidence
	& Vol 2, 5.6-5.7
	&  LO 9.4-9.5
\stepcounter{mylnum}\\ \hline 
\ifnum\lessonNum=\value{mylnum} \rowcolor{green!20!white}\fi
	\themylnum	
	& Oblique Incidence			
	& Vol 2, 5.6-5.7
	& LO 9.5
\stepcounter{mylnum}\\ \hline 
\ifnum\lessonNum=\value{mylnum} \rowcolor{green!20!white}\fi
	\themylnum	
	& Oblique Incidence			
	& Vol 2, 5.6-5.7
	& LO 9.5
\stepcounter{mylnum}\\ \hline 
\ifnum\lessonNum=\value{mylnum} \rowcolor{green!20!white}\fi
	\themylnum	
	& Incidence Lab			
	& TBD
	& LO 9.5
\stepcounter{mylnum}\\ \hline 
%\ifnum\lessonNum=\value{mylnum} \rowcolor{green!20!white}\fi
%	\themylnum	
%	& Wrapup
%	& 
%	& 		 					
%\stepcounter{mylnum}\\ \hline 

\end{tabular}			
\end{\myfontsize}
\end{center}

	\end{textblock}
\end{frame}

%%%%%%%%%%%%%%%%%%%%%%%%%%%%%%%%%%%%%%%%%%%%%%%%%%%%%
%%%%%%%%%%%%%%%%%%%%%%%%%%%%%%%%%%%%%%%%%%%%%%%%%%%%%

\def\obja{ I can apply Coloumb's Law to use direct integration to calculate the electric field due to charge density in varoius geometries, including rectangular, cylindrical, and spherical coordinate systems. }
\def\objb{ I can apply the integral and differential forms of Gauss' Law to calculate the electric field due to charge density in various geometries, including rectangular, cylindrical, and spherical coordinate systems. }
\def\objc{ I can calculate the electric field from the scalar electric potential in a variety of geometries.}
\def\objd{ I can apply the electric field boundary conditions to calculate electric fields on both sides of a material boundary. }
\begin{frame}{Ch 5 Learning Objectives}
	\begin{textblock}{14}(0.5,2.5)
	\begin{myblock}[0.95]{}
	\begin{small}
	\begin{enumerate}[{5}.1]
		\ifnum\objNum=1
			\item \alert{ \obja }
		\else 
			\item \obja
		\fi 
		\ifnum\objNum=2
			\item \alert{ \objb }
		\else 
			\item \objb
		\fi 
		\ifnum\objNum=3
			\item \alert{ \objc }
		\else 
			\item \objc
		\fi 
		\ifnum\objNum=4
			\item \alert{ \objd }
		\else 
			\item \objd
		\fi 	
	\end{enumerate}
	\end{small}
	\end{myblock}
	\end{textblock}
\end{frame}

%%%%%%%%%%%%%%%%%%%%%%%%%%%%%%%%%%%%%%%%%%%%%%%%%%%%%%
%%%%%%%%%%%%%%%%%%%%%%%%%%%%%%%%%%%%%%%%%%%%%%%%%%%%%%
\section{Scalar Potential}
\begin{frame}{Scalar Potential: Line Charge}
\begin{textblock}{7}(0.5,4)
\begin{tikzpicture}[scale=0.7]
		\def\qonex{0}
		\def\qoney{3}
		\def\qtwox{0}
		\def\qtwoy{-2}
		\def\qthreex{4}
		\def\qthreey{0}
		\begin{axis}[ 
		domain=-5:5,
%		samples=200, 
%		xmin=-1, 
		xmax=2.5, 
		ymax=3.5,
		ymin=-3.5,
%		title={Single Square Pulse, $\Pi (\frac{t}{\tau})$},
%		scaled x ticks=base 10:6,
%		xtick scale label code/.code={},
%		minor x tick num = 5,
%		minor y tick num = 5,
		ytick={-3,3},
		yticklabels={-a,a},
		xtick=\empty,
%		xticklabels={$\rho$},
		xlabel={$\rho$},
		ylabel={$z$},
		axis y line = middle,
		axis x line = middle,
		x label style={xshift=0.5cm},
		x label style={yshift=-0.25cm},
		y label style ={xshift=-0.75cm},
		y label style = {yshift =0.5cm},
		yticklabel shift=0.5cm,
		clip=false,
		]
%		\node[cylinder, shape border rotate=90, draw, minimum height=6.5cm, minimum width=0.5cm] at (0,0) {};
		\addplot[ultra thick, red] coordinates {(0,-3)(0,3)};
		\node[red] at (-0.5,0) {$\rho_\ell$};
		\addplot[blue, dashed, ->] (0,1.5) -- (1.5,0) node[midway, above, blue]{r};
		
		\end{axis}
\end{tikzpicture}
\end{textblock}
\begin{textblock}{7}(8,2.25)

	\visible<0|handout:1>{
	\alert{
	\begin{small}
	\eq{
        	&V = \int \dfrac{dQ}{4\pi\varepsilon r} \\
        	&dQ = \rho_\ell d\ell = \rho_\ell dz \\
        	& r = \sqrt{z^2 + \rho^2} \\
        	&\implies  \int \dfrac{\rho_\ell dz}{4\pi\varepsilon \sqrt{z^2 + \rho^2} } \\
        	&= \dfrac{\rho_\ell}{4\pi \varepsilon} \int \dfrac{1}{\sqrt{z^2 + \rho^2}} dz \\
        	&=\dfrac{2\rho_\ell}{4\pi \varepsilon} \lb \text{ln} \lp z + \sqrt{z^2 + \rho^2} \rp \rb\Bigg|_0^a \\
        	&V=\dfrac{\rho_\ell}{2\pi\varepsilon}\lb \text{ln} \lp \dfrac{a + \sqrt{a^2 + \rho^2}}{\rho} \rp \rb
	} % end eq
	\end{small}
	} %end alert
	} % end visible
	
	
\end{textblock}
\end{frame}

%%%%%%%%%%%%%%%%%%%%%%%%%%%%%%%%%%%%%%%%%%%%%%%%%%%%%%
%%%%%%%%%%%%%%%%%%%%%%%%%%%%%%%%%%%%%%%%%%%%%%%%%%%%%%

\begin{frame}{Scalar Potential: Line Charge}
\begin{textblock}{7}(0.5,4)
\begin{tikzpicture}[scale=0.7]
		\def\qonex{0}
		\def\qoney{3}
		\def\qtwox{0}
		\def\qtwoy{-2}
		\def\qthreex{4}
		\def\qthreey{0}
		\begin{axis}[ 
		domain=-5:5,
%		samples=200, 
%		xmin=-1, 
		xmax=2.5, 
		ymax=3.5,
		ymin=-3.5,
%		title={Single Square Pulse, $\Pi (\frac{t}{\tau})$},
%		scaled x ticks=base 10:6,
%		xtick scale label code/.code={},
%		minor x tick num = 5,
%		minor y tick num = 5,
		ytick={-3,3},
		yticklabels={-a,a},
		xtick=\empty,
%		xticklabels={$\rho$},
		xlabel={$\rho$},
		ylabel={$z$},
		axis y line = middle,
		axis x line = middle,
		x label style={xshift=0.5cm},
		x label style={yshift=-0.25cm},
		y label style ={xshift=-0.75cm},
		y label style = {yshift =0.5cm},
		yticklabel shift=0.5cm,
		clip=false,
		]
%		\node[cylinder, shape border rotate=90, draw, minimum height=6.5cm, minimum width=0.5cm] at (0,0) {};
		\addplot[ultra thick, red] coordinates {(0,-3)(0,3)};
		\node[red] at (-0.5,0) {$\rho_L$};
		\addplot[blue, dashed, ->] (0,1.5) -- (1.5,0) node[midway, above, blue]{r};
		
		\end{axis}
\end{tikzpicture}
\end{textblock}
\begin{textblock}{7}(8.5,2.75)

	\visible<0|handout:1>{
	\alert{
	\begin{small}
	\eq{
        	&V=\dfrac{\rho_\ell}{2\pi\varepsilon}\lb \text{ln} \lp \dfrac{a + \sqrt{a^2 + \rho^2}}{\rho} \rp \rb \\
        	&\implies\vec{E} = -\nabla V = -\rhohat \partd{V}{\rho} \\
        	&\text{for } a >> \rho \text{ :} \\
        	&\implies \vec{E} = -\rhohat \partd{}{\rho}\dfrac{\rho_\ell}{2\pi\varepsilon}\lb \text{ln} \lp \dfrac{2a }{\rho} \rp \rb \\
        	& = -\rhohat \partd{}{\rho}\dfrac{\rho_\ell}{2\pi\varepsilon}\lb \text{ln} \lp 2a \rp - \text{ln}\lp \rho\rp \rb \\
        	&\implies \vec{E} = \rhohat \dfrac{\rho_L}{2\pi \varepsilon \rho}     	
	} % end eq
	\end{small}
	} %end alert
	} % end visible
	
\end{textblock}
\end{frame}

